%--
\chapter{函数的渐近增长率}

%----------------------------------------
\section*{章节注释}

%----------
\subsubsection{微积分严格化的历史}

$\epsilon-\delta$语言的背后，是微积分严格化的艰苦历程。

牛顿和莱布尼茨虽然把微积分系统化，但是它还是不够严谨。
可是当微积分被成功地用来解决许多问题，却使得十八世纪的数学家偏向其应用，而甚少致力于其严谨性。
一些数学家，包括科林·麦克劳林，试图证明使用无穷小是可靠的做法，但直到150多年之后才得以成功。

奥古斯丁·路易·柯西和卡尔·魏尔斯特拉斯的工作，实现了对无穷小的记号的回避。
微分和积分的基础终于被打下了。
在柯西的著作中，可以看到一系列的基础进路尝试，包括通过无穷小来对连续进行定义，和在微分定义中一个不太精确的函数极限原型。
而在魏尔斯特拉斯的著述中，对极限概念作了形式化，回避了无穷小的使用。
继魏尔斯特拉斯之后，微积分就常以极限作为基础，而非无穷小了。